\documentclass[a4paper,11pt,twoside,openright]{book}

\usepackage[spanish]{babel}
\usepackage[utf8]{inputenc}

\usepackage{eurosym}
\usepackage{fancybox}
\usepackage{multicol}
\usepackage{amsmath}
\usepackage{subcaption}
\usepackage[nottoc,numbib]{tocbibind}
\usepackage[pdftex,usenames,dvipsnames]{color}

% compilar con pdflatex.
\usepackage[pdftex]{graphicx}


\usepackage{setspace}
\onehalfspacing
\usepackage{geometry}
\geometry{
  inner=3.5cm,
  outer=2.5cm,
  bottom=3.5cm,
  top=3.5cm}

\usepackage{fancyhdr}
\pagestyle{fancy}
\fancyhf{}
\fancyhf[HR]{\thepage}
\fancyhf[HL]{\nouppercase\rightmark}

\usepackage{booktabs}
\usepackage{hyperref}
\hypersetup{
    colorlinks,
    linkcolor={red!50!black},
    citecolor={blue!50!black},
    urlcolor={blue!80!black}
}
\usepackage{rotating}
\usepackage{multicol}
\usepackage{multirow}
\usepackage{pgfgantt}
\usepackage{enumitem}
\usepackage{color}
\usepackage{xcolor}
\usepackage{caption}
\DeclareCaptionFont{white}{\color{white}}
\DeclareCaptionFormat{listing}{\colorbox{gray}{\parbox{\textwidth}{#1#2#3}}}
\captionsetup[lstlisting]{format=listing,labelfont=white,textfont=white}

 \usepackage{listings}
 \usepackage{courier}
 \lstset{
         basicstyle=\footnotesize\ttfamily, % Standardschrift
         numberstyle=\tiny,          % Stil der Zeilennummern
         numbersep=5pt,              % Abstand der Nummern zum Text
         tabsize=2,                  % Groesse von Tabs
         extendedchars=true,         %
         breaklines=true,            % Zeilen werden Umgebrochen
         keywordstyle=\textbf,
         frame=b,
         stringstyle=\textit, % Farbe der String
         showspaces=false,           % Leerzeichen anzeigen ?
         showtabs=false,             % Tabs anzeigen ?
         xleftmargin=17pt,
         framexleftmargin=17pt,
         framexrightmargin=5pt,
         framexbottommargin=4pt,
         %backgroundcolor=\color{lightgray},
         showstringspaces=false      % Leerzeichen in Strings anzeigen ?
 }

 \lstset{literate=
  {á}{{\'a}}1 {é}{{\'e}}1 {í}{{\'i}}1 {ó}{{\'o}}1 {ú}{{\'u}}1
  {Á}{{\'A}}1 {É}{{\'E}}1 {Í}{{\'I}}1 {Ó}{{\'O}}1 {Ú}{{\'U}}1
  {à}{{\`a}}1 {è}{{\`e}}1 {ì}{{\`i}}1 {ò}{{\`o}}1 {ù}{{\`u}}1
  {À}{{\`A}}1 {È}{{\'E}}1 {Ì}{{\`I}}1 {Ò}{{\`O}}1 {Ù}{{\`U}}1
  {ä}{{\"a}}1 {ë}{{\"e}}1 {ï}{{\"i}}1 {ö}{{\"o}}1 {ü}{{\"u}}1
  {Ä}{{\"A}}1 {Ë}{{\"E}}1 {Ï}{{\"I}}1 {Ö}{{\"O}}1 {Ü}{{\"U}}1
  {â}{{\^a}}1 {ê}{{\^e}}1 {î}{{\^i}}1 {ô}{{\^o}}1 {û}{{\^u}}1
  {Â}{{\^A}}1 {Ê}{{\^E}}1 {Î}{{\^I}}1 {Ô}{{\^O}}1 {Û}{{\^U}}1
  {œ}{{\oe}}1 {Œ}{{\OE}}1 {æ}{{\ae}}1 {Æ}{{\AE}}1 {ß}{{\ss}}1
  {ç}{{\c c}}1 {Ç}{{\c C}}1 {ø}{{\o}}1 {å}{{\r a}}1 {Å}{{\r A}}1
  {€}{{\EUR}}1 {£}{{\pounds}}1
 }

\parskip=6pt

\parindent=10pt

\setcounter{secnumdepth}{3}

\title{Desarrollo de un videojuego roguelike para invidentes aplicando técnicas de Procesamiento del Lenguaje Natural.}
\author{Darío Penas Sabín}
\date{\today}


\renewcommand\lstlistingname{Fragmento de Código}
\renewcommand\lstlistlistingname{Fragmentos de Código}

\newenvironment{bottompar}{\par\vspace*{\fill}}{\clearpage}

\begin{document}
        %
% Portada.
%

% Nota: Sería más cómodo emplear el comando \maketitle que genera una portada de forma automática, pero
% no incluye toda la información que es necesario incluir en la memoria de un proyecto de fin de carrera
% de la Facultad de Informática de A Coruña.
%

\begin{titlepage}

	\begin{center}

		% Logotipo de la universidad.
		\includegraphics[width=6cm]{./eps/logo_udc.eps}
		\vspace{2cm}

		{\Large{\textbf{Facultade de Informática da Universidade de A Coruña}}}
		\\
		{\it \large{\textbf{Computación}}}
		\vspace{1cm}

		% Indicamos el nombre de la titulación oficial que hemos cursado con tanto esfuerzo.
		{\large PROYECTO DE FIN DE CARRERA\\Ingeniería Informática}
		\vspace{1cm}

		% Título
		\textbf{\Large Desarrollo de un videojuego roguelike para invidentes aplicando técnicas de Procesamiento del Lenguaje Natural.}
		\vspace{7cm}
	\end{center}

	\begin{flushright}
		\begin{tabular}{ll}
			\large{\textbf{Alumno:}}	&
			\large{Darío Penas Sabín} \\

			% Nombre del director/tutor del proyecto.
			\large{\textbf{Director:}}	&
			\large{Jesús Vilares Ferro} \\

			% Nombre del director/tutor del proyecto.
			\large{\textbf{Director:}}	&
			\large{Carlos Gómez Rodríguez} \\

			% Fecha.
			\large{\textbf{Fecha:}}	&
			\large{CHANGE: 15 de junio de 2016} \\
		\end{tabular}
	\end{flushright}

\end{titlepage}


        \frontmatter

        \thispagestyle{empty}

        \section*{Resumen:}

La industria del entretenimiento digital ha crecido inmensamente en los últimos años, llegando a alcanzar números de ventas jamás vistos anteriormente.
Parte de la razón de este crecimiento viene dada por una mejora radical en el aspecto visual, necesaria para que el jugador se sienta inmerso en la aventura que se le está planteando.
Estas mejoras, sin embargo, dejan de lado a muchos jugadores que, por diferentes motivos, no son capaces de apreciar el contenido visual que se les ofrece o tienen problemas para ello, haciendo imposible que disfruten del contenido ofrecido.

Este proyecto consiste en la creación de un videojuego de género \textit{roguelike} para invidentes que, desde un principio, parte de la idea de generar contenido específicamente diseñado para que pueda ser jugado por todo el mundo, haciendo énfasis en ofrecer al jugador una diversa cantidad de frases que describan lo que está sucediendo en su alrededor y que serán generadas automáticamente en base a una serie de gramáticas y diccionarios, empleando para ello técnicas de Procesamiento del Lenguaje Natural.

        \thispagestyle{empty}

        \section*{Lista de palabras clave:}

\begin{itemize}
  \item Tiflotecnología
  \item Procesamiento del Lenguaje Natural
  \item Accesibilidad
  \item Entretenimiento Digital
  \item \textit{Roguelike}
  \item Java
  \item \textit{Open Source}
\end{itemize}

        \thispagestyle{empty}

        %
% Agradecimientos
%

\section*{Agradecimientos}

Me gustaría aprovechar este pequeño espacio para agredecer a todas las personas que me ayudaron, en diferentes niveles y formas, a elaborar y terminar este proyecto.

Gracias a todos aquellos que dedicaron parte de su tiempo a resolverme dudas sobre tiflotecnología y accesibilidad (sobre todo a los usuarios de \textit{Reddit} \textit{ais523} y \textit{fastfinge}), así como a todos los profesores que tuve durante los cinco años de carrera y que lograron que llegue al punto en el que me encuentro profesionalmente. En especial me gustaría agredecer a mis directores, Jesús y Carlos, el esfuerzo extra que han hecho durante este tiempo en el que he estado trabajando en el extranjero para mantenerme al tanto y resolverme dudas a distancia cuando estaba perdido.

Tampoco puedo olvidarme de todos los amigos que he hecho durante todos estos años y con los que he pasado tantas horas delante y fuera de la pantalla. Adrián (insaciable compañero de prácticas que me tuvo que aguantar durante tantos años), Daniel, Vanesa, Santi, Chava, Mónica, Brais, María, Xacobe y otros muchos. ¡Gracias por vuestra ayuda y amistad!

Por último, gracias a todas aquellas personas (especialmente a mis padres y a mi novia Diantha) que me ayudaron y presionaron a que terminara el proyecto cuanto antes y no lo dejara pasar.

¡Gracias a todos por todos estos maravillosos años!

\begin{flushright}
  Darío Penas Sabín \\
  Amsterdam, \today
\end{flushright}


        \thispagestyle{empty}

        \mainmatter
        \chapter[Introdución]{Introdución}

En este capítulo introductorio se explicarán los aspectos necesarios para entender lo más importante del proyecto, la motivación para realización del mismo y un breve resumen del resto de capítulos que forman parte de la memoria.

\section{Videojuegos y personas invidentes}
La mayor parte de los videojuegos comerciales no tienen en cuenta a muchas minorías de la sociedad. Haciendo una pequeña búsqueda online pueden encontrarse miles de personas quejándose de \textit{first person shooters} que tienen un \textit{FOV}\footnote{\textit{Field of view}, campo de visión. Extensión de mundo observable en un momento dado} limitado, causándoles mareos al poco rato; daltónicos protestando que diferentes juegos (como por ejemplo The Witness\footnote{Juego de puzzles en primera persona: \url{http://the-witness.net/}}), basan buena parte de su mecánica en que el jugador sea capaz de distinguir diferentes colores; zurdos que tienen que acomodarse a ciertos controles a no existir una opción para cambiarlos; invidentes que no pueden disfrutar de prácticamente ninguno de los títulos que se encuentran en el mercado, etc.

En este proyecto nuestro objetivo es crear un videojuego desde cero que tenga en cuenta todo tipo de minorías, tomando especial relevancia el desarrollo para invidentes y centrándose en los aspectos que sean relevantes para ellos; descripciones que sean fácilmente reproducibles que cambien automáticamente y fácil expansión del título, tanto en características generales como en idiomas, gramáticas o palabras empleadas en el mismo.

\subsection{Visita a la \textit{ONCE}}
A principios de 2015 Jesús y yo fuimos a un taller de la ONCE donde una persona invidente nos habló sobre la tiflotecnología, nos mostró la forma en la que usaban ordenadores y móviles, incluso para leer código y los errores, muchos de ellos fácilmente solventables, que se cometían día a día en temas de accesibilidad. Esta visita nos abrió los ojos y gracias a ella fuimos capaces de detectar posibles mejoras que hacer al proyecto y los fallos que no deberíamos de cometer en su diseño. También se ofrecieron a probar el proyecto una vez estuviera listo, pero este será un tema que trataremos en las próximas secciones.

\section{Motivación}
El entretenimiento digital siempre ha sido parte de mi vida y una de las razones por las que desde pequeño estuve interesado en la informática y motivo por el que, finalmente, acabé estudiando esta carrera. Del mismo modo, y habiéndome relacionado con bastante gente con toda clase de necesidades especiales durante un gran periodo de tiempo, mi interés por la tiflotecnología y las limitaciones que la tecnología ofrece a millones de personas no hizo más que crecer año a año.

A pesar de los grandes avances de la industria de los videojuegos y de la gran cantidad de nuevos estudios y proyectos que se lanzan anualmente, resulta muy complicado poder dedicarse profesionalmente a cualquiera de estas dos cosas (y no digamos ambas a la vez) por lo que en mi futuro profesional no he sido capaz, al menos de momento, de cumplir mi sueño de trabajar en lo que más me apasiona. Por este motivo, cuando Jesús me comentó que él y Carlos llevaban un tiempo con este proyecto disponible, no dudé en un instante en aceptarlo y ponerme manos a la obra.

Este videojuego se ha desarrollado durante el periodo de un par de años en los que se incluyen muchos cambios en mi vida, tales como mi estancia permanente en Holanda hace ya casi dos años y mis primeros pasos en el mundo laboral. Cada día que pasa me alegro más de tener un proyecto como éste, dado que sin la pasión e interés por el mismo estoy seguro de que jamás lo habría terminado.

\section{Estructura de la memoria}

La memoria está formada por ocho capítulos en los que se explican los pasos tomados a la hora de crear el proyecto.

\paragraph*{Capítulo 1. Introducción.}
Se explicarán, de manera general y resumida, en qué consiste el proyecto, la motivación del mismo y la estructura que tendrá la memoria.

\paragraph*{Capítulo 2. Estado del arte.}
En este apartado hablaremos de otros proyectos similares, de la situación actual de la industria sobre el problema que tratamos en este proyecto y del diferente software que es utilizado en relación con la tiflotecnología.

\paragraph*{Capítulo 3. Fundamentos Tecnológicos.}
Citaremos y hablaremos sobre las herramientas y bibliotecas empleadas durante la elaboración del proyecto.

\paragraph*{Capítulo 4. Metodología.}
Se detallarán las prácticas y metodologías de desarrollo empleadas para la realización del videojuego y la razón de su uso.

\paragraph*{Capítulo 5. Planificación y Seguimento.}
Detallaremos la planificación y seguimiento usados en cada una de las etapas del proyecto. 

\paragraph*{Capítulo 6. Análisis de requisitos.}
Se comentará el análisis de requisitos para este proyecto y se explicará en detalle cada uno de los mismos.

\paragraph*{Capítulo 7. Diseño e implementación.}
En este capítulo explicaremos los detalles del diseño y de la implementación de ciertas partes del programa que consideramos más relevantes.

\paragraph*{Capítulo 8. Recepción, conclusión y trabajo futuro.}
En este apartado mostraremos los comentarios obtenidos por la comunidad durante el desarrollo del juego, se relatarán las conclusiones obtenidas y se detallarán los posibles cambios, mejoras y extras que se podrán tener en cuenta en el futuro para su mejora.

        \chapter{Fundamentos Tecnológicos}

En este capítulo hablaremos sobre los fundamentos tecnológicos que vamos a usar en este proyecto y, si cabe, la razón por la que fueron elegidas. En primer lugar citaremos las herramientas que hemos usado y, en segundo lugar, las bibliotecas que hemos decidido utilizar en el programa en sí.

\section{Herramientas empleadas}

\paragraph{Java} Lenguaje de programación orientado a objetos cuya primera aparición fue en 1995. Es uno de los lenguajes de programación más utilizados en la industria y una de sus principales características es que es multiplataforma, es decir, puede ser ejecutado en cualquier sistema operativo que tenga la \textit{Java Virtual Machine} instalada sin necesidad de realizar cambios en el código (WORA\footnote{\textit{Write once, run anywhere}. Eslogan creado por Sun Microsystems para mostrar los beneficios de la multiplataforma}). Esta ventaja es esencial en nuestro caso, dado que la mayoría de las personas que pueden estar interesadas en el proyecto usan una gran variedad de sistemas operativos.

 \paragraph{Eclipse} Es un IDE\footnote{\textit{Integrated Development Environment}. Entorno de desarrollo integrado} usado para escribir código en múltiples idiomas. También incluye una serie de \textit{plugins} que facilitan y automatizan muchas de las labores a realizar como el uso de sistema de controles, ejecución de código y tests, herramientas de debug, autocompletado de código, etc.

 \paragraph{Git} Sistema de control de versiones distribuido introducido en 2005 y desarrollado principalmente por Linus Torvalds. Es el control de versiones referencia en la mayoría de empresas y proyectos de software libre gracias a su rapidez y, al ser distribuida, permite trabajar y realizar \textit{commits} del código sin necesidad de conexión a internet.

 \paragraph{GitHub} Plataforma de desarrollo colaborativo usada para alojar proyectos usando el sistema de control de versiones Git. La mayoría de proyectos de código abierto lo usan, dado que es gratuito, aunque también tiene la opción de almacenar el código de forma privada tras, previamente, realizar un pago.

 \paragraph{Listas de Correo} Las listas de correo son un método de comunicación muy usado por diferentes comunidades, especialmente en el desarrollo de software, que ayudan a los usuarios que participan en ellas a enviar correos a múltiples personas que lo deseen de forma anónima y, al mismo tiempo, tener un historial de las respuestas dadas por los mismos. En nuestro caso la hemos usado para comunicarnos con un grupo de usuarios y desarrolladores de videojuegos para invidentes.

 \paragraph{Reddit} Web creada en 2005 y que actualmente se encuentra en el top 50 de las más visitadas del mundo. Cuenta con una comunidad gigante que está dividida en muchísimos subgrupos dependiendo del tema a tratar. La hemos usado como una herramienta de feedback. Especialmente los \textit{subreddits} de daltónicos\url{https://www.reddit.com/r/ColorBlind/} y gente ciega\url{https://www.reddit.com/r/blind/}. 

\paragraph{Dia} Aplicación informática que permite la creación de todo tipo de diagramas. En nuestro caso lo hemos usado para crear los diagramas UML que se encuentran en esta memoria.

\paragraph{Gantt Project} Programa de software libre diseñado para la creación de diagramas de Gantt.

\paragraph{JSON} JavaScript Object Notation. Es un formato muy usado en APIs para intercambio de datos, similar a XML. En nuestro caso lo usamos para definir las gramáticas y diccionarios de nuestro proyecto, dado que es muy sencillo de leer y especificar. Hay numerosas bibliotecas que nos permiten analizar y trabajar con este formato en Java. La que nosotros usamos es \textit{Gson}.

\paragraph{\LaTeX} Sistema de composición de textos altamente usado por la mayoría de textos científicos dada la facilidad de su composición, simpleza, alta calidad y herramientas que a ayudan a la creación de fórmulas, inserción de imágenes y muchos otros elementos. Muy modificable. Es el sistema que hemos usado para la creación de este documento.

\paragraph{NVDA} Lector de pantalla de código libre para Windows. Orca es, en cierta medida, su equivalente en Linux.

\section{Bibliotecas empregadas}

\paragraph{Gson} Biblioteca usada para transformar archivos JSON a objetos de Java y viceversa.

\paragraph{JCurses} JCurses\footnote{\textit{The Java Curses Library}} es una biblioteca para el desarrollo de aplicaciones de terminal para JAVA. Es similar a AWT\footnote{\textit{Abstract Window Toolkit.} Kit de herramientas de interfaz de usuario de la plataforma original de Java}, pero basada en el sistema de ventanas Curses de UNIX.

\paragraph{Libjcsi} Biblioteca de representación gráfica que trabaja sobre JCurses y simplifica la tarea de representar y refrescar elementos del terminal.


\end{document}

